\section{Beyond Mediocrity}

\begin{quotex}
The Mage thinks and wills; he loves nothing with desire; he rejects nothing in passion. 
\flright{\textsc{Eliphas Levi}}
\end{quotex}

Although \textbf{Dostoyevsky} wrote the The Idiot\footnote{\url{https://en.wikipedia.org/wiki/The\_Idiot}} in 1868, his depiction of human characters is so accurate, it could have been written yesterday afternoon. There are only a limited number of types of people; certain archetypes are incarnated, and prototypical situations are repeated ad nauseum in history.

\begin{quotex}
There are some people about whom it is difficult to say anything which would describe them immediately and fully in their most typical characteristic aspects; these are the people who are usually called ``ordinary'' and accounted as the ``majority''.
\end{quotex}


Of course, political parties praise the ordinary and the ``silent'' majority. And rightfully so. In a well- organized society, the ordinary provide for the needs and servicing of economic life. When there is harmony in a hierarchic arrangement, the ordinary may be inarticulate, but they have an animal faith and sound instincts, particularly if they are insulated from the ideologies of the pseudo-intelligentsia. A leader, who knows what the ordinary believe, is able to articulate and make conscious that faith and instincts, that is, he awakens what had previously lain dormant in the masses.

\begin{quotex}
When, for example, the very nature of certain ordinary persons consists precisely in their perpetual and unvarying ordinariness, or better still, when in spite of their most strenuous efforts to lift themselves out of the rut of ordinariness and routine, then such persons acquire a certain character of their own --- the typical character of mediocrity which refuse to remain what it is and desires at all costs to become original and independent, without having the slightest capacity for independence.
\end{quotex}


In an attempt to find their own voice, they will latch on to a religion, ideology or naïve faith in science. Yet, here they remain dependent on the authority of the leaders in each of those fields.

\begin{quotex}
There are a great number of such people in the world, far more than it appears. Like all people, they may be divided into two categories: some are mentally limited, others, much cleverer. The first are happier. For the ordinary person of limited intelligence nothing is easier than to imagine himself an exceptional and original person and to take delight in this delusion with no misgivings. It has been enough for some of our young ladies to cut their hair short, put on blue eyeglasses, and call themselves Nihilists for them to persuade themselves that, in putting on their spectacles, they immediately acquired convictions of their own. It is enough for a man to feel in his heart a droplet of humanitarian and benevolent emotion to be immediately persuaded that no one feels so deeply as he and that he stands at the very vanguard of civilization. It is enough for another to pick up some thought he has heard, or to read a page at random somewhere, to believe at once that it was ``his own idea,'' engendered in his own brain.
\end{quotex}


These types of characters are amazingly still with us. The pseudo-rebellion just based on changing appearance; the ``highly-evolved'' free spirit who cares for humanity, animals, and nature. The autodidact who latches onto an idea isolated from its larger context. I don't want to heap unnecessary scorn on those of a sincere nature who oppose the status quo. The problem is that it is all done horizontally, to oppose one conformist group with another. It is all emotionally driven. Only a very few will move in a vertical direction, that is they will transcend the ordinary, not oppose it. These are the Initiates. The emotional and inconclusive battles between ideologies, or religious and political beliefs, are no longer of interest to them.


\flrightit{Posted on 2010-11-09 by Cologero}

\begin{center}* * *\end{center}

\begin{footnotesize}
\begin{sffamily}
\texttt{kadambari on 2010-11-10 at 12:47 said:}

There is nothing quite like reading the old Russian masters of literature. I owe my exposure to Russian literature to a Russian family friend who gave me copies of Russian literature translated into English, telling my parents what I was reading at the time did not constitute ``literature''. The Idiot is perhaps one of the most beautiful novels I have ever read.

Can you give your examples of ``initiates'' in history, who exhibits vertical developments? Even people like Gandhi appear of limited vision when you really read his autobiography (which was my experience), and as wanting to impose their limited views on others. This is also true for people like Mother Teresa. The only ``initiate'' I can think of is the Sakyamuni, and enlightened emperors and rulers to a lesser extent. The great artists and philosophers in history would also be under this category...

\hfill

\texttt{kadambari on 2010-11-11 at 09:58 said:}

above

who ``exhibit''

Also the problem has always to bring the contemplative man (contemplative in the sense of evolved intellectually and spiritually) back to the polis, how can the man depicted in book ten of Aristotle's Ethics be persuaded to come back to the polis and take part in organizing its affairs? Even the great Aristotle does not offer answers and leaves this unresolved. We have not moved any closer to solving this problem, how can a man who is ``beyond'' politics be expected to take part and be active in a system through which he clearly sees and which he does not necessarily believe in? Only when he believes that not doing anything would result in a worse state of affairs (solution of Plato).

\hfill

\texttt{Cologero on 2010-11-11 at 11:07 said:}

Yes, good point. That is why we emphasize sacred magic. Thought leads to action; the truth of an idea (essence) must be reflected in existence ... beyond thought, there is Will. We have emphasized that repeatedly. The solution to the world does not lie in thought. The true man is beyond political disputes, not politics per se, but he is not torn by the yes and the no, the debate, since he understands the one. He is the Emperor.

\hfill

\texttt{Graham on 2010-11-22 at 21:07 said:}

kadambari,

What are the best books on Persian history and culture?

\hfill

\texttt{kadambari on 2010-12-02 at 13:46 said:}

@Graham

I am puzzled by your question. I am not Persian, and my interest in Persian history is limited to the history of pre-Islamic Persia and the few centuries that comprised the golden age of Islam, which was Persian dominated. I am still reading and learning in this respect.

I think for beginners, Will Durant's ``Our Oriental Heritage'' has a short section on Persia, and has a good bibliography. I started learning world history in high school with Durant's Story of Civilization which is a good introduction to world civilizations; people do not write history in this manner these days; the writing of history has mostly become specialized and focused on minutiae, with the large picture ignored. Toynbee's A Study of History is also a good read in this respect to obtain a large picture.

The history of Persia is complex and is composed of many layers. Modern day Iranians are not wholly the same peoples as the ancient Persians, as they were conquered by the Turks (not to be confused with Ottoman Turks) and Arabs (hence, continuity in culture has been disrupted, resulting in the two conflicting cultural identities of modern day Iranians ---that borne out of ancient Persia and that stemming from the Islam); however, they retain their language which is a lot compared to others who lost their language completely to Arabic (Lebanese, Syrians). My understanding is that they are the most intelligent amongst the peoples in the Islamic world.

The dominant culture in Iran is Shia Islam today which it adopted in the 1501 (the Safavids declared Shiism as the State religion).

High Persian culture ended when the Arabs conquered Persia. These people, unlettered, tribal and used to living in the harsh desert where there was a constant fight for survival, did not possess much else other than an oral tradition. It took almost a century for the Persians to be converted by the sword although they put up a brave resistance. This conquest of Persia is one of the most tragic tales of world history (the other being the Islamic conquest of India). Although the religion of the Persians changed after this tragic event, the Persians in a sense conquered the Arabs by giving their culture and learning to Islam. Most of the so-called Islamic men of science during the Golden Age of Islam were Persian and the Golden Age of Islam (8-11 centuries) which set the standard for Islamic culture was Persian dominated. For example, the names of the Islamic men of science that people are familiar with such as Avicenna, Khayyam, Razi, Khwarizmi (algebra), Naubakht, Ibn Muquaffa, Banu Moussa, Qazwini, Majusi and countless others were Persian in culture. Many of the writings were composed in Arabic because people were forced to write in Arabic until Persians recovered their Persian language through Firdusi (epic poet of Persia).

The Zoroastrians have become a negligible community in Iran and have mostly emigrated to the West. Some of the Persian aristocracy that did not want to convert to Islam emigrated to India in ships and landed on the Western Coast of India (Gujarat). The make a pact with the local ruler who let them stay as long as they did not disrupt the local culture. The Parsis were poor farmers and no one cared much about them. However, under the British, they became the business henchmen of the British. Many of them made money through the opium trade, which the British devised to balance the trade deficit with China. Like today, China sold goods like tea and silk to the West but did not like to buy WEstern goods. So via the opium trade, 1/3 of the Chinese (roughly 125 million became addicted to opium), reversing the trade deficit. It is interesting that a recent Wall Street journal article was called ``China's drug resistant budget surplus''! The Parsis in India are a highly educated class today and are industrialists (such as the famous TATA), they are more a mercantile community than an intellectual community.

Now why after such a rich tradition of thinking and creativity does Persia slide after the 1200's? My take on it is that the essential barren nature of Islam sets in and the culture becomes more Islamized than Persian, although Persian creativity lasts for a few hundred years after the conquest. The rest of their history is a black hole in terms of culture and creativity afterwards. It is easy to see why the West gets ahead of the East in terms of science after this period.

A comparison can be made to India. Indian history is pretty much a black hole as well after the Islamic conquests, there is no spectacular creativity in the culture anymore and native learning and creativity dries up and withers after the destruction. You have some gaudy structures like the Taj Mahal (hardly an architectural wonder, unless you call a building made of marble and gems an architectural wonder!). The native architecture stressed harmony with nature, and the native architecture blends in with nature and was harmonious with the surroundings unlike the gaudy structures such as the Taj Mahal, built on a different conception of architecture. There are some fine arts like love poetry in the Islamic dominated court but that is it in terms of culture. Islam occupying West Asia to Northern India, destroys the exchange of ideas between the East and West (the West after Islam knows of India only thorough the Arabs who take credit for Indian achievements, and Columbus has to sail the world to rediscover India, although the Greeks were frequent visitors there). This intellectual isolation under Islam ends with the British, who although as colonizers have the sole aim of making a buck from the region, bring in new ideas, establish universities to train bureaucrats, and end the intellectual black hole under Islam. At least they do not loot temples for their wealth, but exploit in a different manner.

Amir Mehdi Badi who came from an old Iranian family was a prolific scholar on ancient Persia. Many of his books are in French though.

Hope this helps a bit.

\hfill

\texttt{kadambari on 2010-12-02 at 16:27 said:}

@Graham

Sorry for the small typos such as make for ``made'', the Islam instead of ``Islam''

and so on. I find it hard to edit on this site before pressing submit... 

\hfill

\texttt{kadambari on 2010-12-03 at 09:20 said:}

I forgot to add that the pre-Islamic history of Persia involves a great deal of reconstruction from Greek and Roman sources. The holy books of the Zoroastrians were destroyed after the Islamic invasions and what what remains today of the Gathas preserved by the Parsis is but a small portion of what existed... I think understanding the patterns of thinking of pre-Islamic Persia, classical India before Islam, classical Greece and Rome provide clues to how people thought and percieved the world before being influenced by the desert cults from the Mid-East... 

\hfill

\texttt{kadambari on 2010-12-09 at 17:39 said:}

@Graham

On fravahr are some articles on Amir Mehdi Badi's writings... My understanding is that the Persians and Greeks back then had more in common than Persians and Greeks today. And later, if the Persians and Romans had not exhausted each other fighting, the world would have been less prone to the damage done by Mid-Eastern dogmatisms... and the nobler vision of mankind might have prevailed... This is my general sense of history, but still need to do a great deal of study and research to undertand this better... 

\hfill

\texttt{kadambari on 2010-12-09 at 17:48 said:}

the link for the above:\\
\url{http://www.fravahr.org/spip.php?rubrique11}


\url{http://www.fravahr.org/spip.php?rubrique11}


\hfill

\texttt{kadambari on 2010-12-09 at 18:12 said:}

The second link was wrong.

\url{http://www.fravahr.org/spip.php?article16}


I meant to say the Romans became Christian and weakened the Persians were conquered by Arabs and changed in their culture; it was easy for the Arabs to conquer the Persians who had been exhausted by fighting the Romans... 

I have always wondered if the backwardness of Greeks to the other Europeans in modern times has to do with them being under the Turks?

The Greek commentators on India are interesting. India is known for honesty, lack of materialism despite plenty (back then) by the Greeks. After centuries of Islamic tyranny and extortion, India is known by famines, for thugish culture, snake charmers and what not... A look at Persian before and after is also an eye opener in this respect... 

\hfill

\texttt{kadambari on 2010-12-09 at 18:14 said:}

error above-a look at Persian culture before and after the foreign conquests, is also an eye opener.

\hfill

\texttt{Blind Dog on 2016-11-14 at 01:28 said:}

Kadambari you are wrong, Islam is great, it was anglo-saxons and losers who ruined India, and where are your proofs that Muslims destroyed the Avesta? Maybe Hulagu and Alexander did it. This will be settled on Judgment Day; Al-Tabari alone is not enough evidence for a verdict.

\hfill

\texttt{Ulysses on 2016-11-14 at 01:43 said:}

Kadambari, please update your links.


\hfill

\end{sffamily}
\end{footnotesize}

